\documentclass[11pt,a4paper]{article}
\usepackage[T1]{fontenc}
\usepackage[utf8]{inputenc}
\usepackage{times}
\usepackage[top=2cm,bottom=1.8cm,left=2.3cm,right=2.3cm]{geometry}
\usepackage{microtype}
\usepackage[hidelinks]{hyperref}
\pagestyle{empty}
\setlength{\parindent}{0pt}
\setlength{\parskip}{8pt}

\begin{document}

{\large\bfseries Grover Aruquipa}\par\vspace{-5pt}
{\small Qu\'ebec, QC, Canada \ \textbullet\ \ grover.aruquipa.9@gmail.com \ \textbullet\ \ groveraruquipa.github.io}

\vspace{2pt}
Re: Postdoctoral Research Assistant in Tensegrity Robot Design and Control (vacancy 189502)

Hi Robert,

I was looking for a postdoc, and also for a place to apply for the Branco Weiss Fellowship next year, and when I saw this position in your lab I thought it was the right fit for both. You already gave me good advice on Branco Weiss in September, so I wanted to write to you directly. I would like to apply for this post, and if it works for you, I would also like to prepare a Branco Weiss application with your lab as host.

A bit about what I do. I am finishing my PhD at Universit\'e Laval with Cl\'ement Gosselin and Patrice Lambert, and I defend in January 2027. My thesis is on parallel robots where the mechanism does the work: all motors are on the base, the gripper is driven through the structure, and the robot can rotate what it holds without limit. I have built three of these robots, with 4, 7 and 8 motors. I did the design in SolidWorks, the machining and printing, the wiring, and I wrote the motor control myself, running at 500\,Hz. The smallest one, a coaxial wrist-gripper, turned its tool ten full turns with less than a degree of error while holding the object. Right now I am building a Delta robot with a gripper that feels the grasp force through the motor currents.

On simulation, I use MuJoCo and Simscape to check my robots before I build them, and I have experimented with Isaac Sim, so moving to IsaacLab should be quick. I have not built a tensegrity yet. But the math I use every day is close to it: a singularity in my parallel robots and a soft mode in a tensegrity are the same kind of rank loss in the equilibrium matrix, and the prestress of a tensegrity lives in the null space of that same matrix. So I think I can help with the design and control side from the first week, and I am honestly excited to learn how to build them.

What I like about your work is that the body of the robot is part of how it solves the task, like in your Nature paper on adaptive morphogenesis and in Spikebot. That is where I want my own research to go. One idea I would like to try is using the shape change of a tensegrity not only to move, but also to grasp and turn objects, which links straight to what I did in my PhD.

I also supervise two graduate students on one of my robots, and I led the Laval team that won the ASME Student Mechanisms and Robotics Design Competition in 2025. I have 3 journal papers as first author and 4 more under review.

I could start in early 2027 and would need a UK Skilled Worker visa. My referees are Cl\'ement Gosselin (clement.gosselin@gmc.ulaval.ca) and Andreas M\"uller at JKU Linz (a.mueller@jku.at), who hosted me this year.

Thanks a lot for your time. I would be very happy to talk whenever it suits you.

Best regards,\\
Grover Aruquipa

\end{document}
